% Provenance: copied from the design workspace CambridgePQ/recursion.fast, `report/sections/08-security.tex`, on 2026-10-19. Read-only copy; the source is authoritative.
% 08-security -- written by W2
\section{Security definitions}\label{sec:security}

This section states what a submission proves about its spec verifier, why each definition has the
form it has, what the definitions buy after compilation to a hash-based argument, and how each could
be made vacuous. The definitions are restated locally over Lean core and Mathlib, in 132 code
lines of contract v0.3 (\code{Shapes.lean} 50, \code{Security.lean} 55, \code{Protocol.lean} 27). They are generic in input and output statements and witnesses, so the base track
(output statement \lean{Unit}, output relation \lean{Set.univ}) and the three accumulation
reductions of \cref{sec:recursion} use the same text. An untrusted \code{Bridge/} module proves
that they imply ArkLib's \lean{rbrKnowledgeSoundnessWorstCaseWith} (\lean{RbrKS.toArk})
\cite{arklib2026}, so a certificate means what ArkLib means. The converse, which would let agents
develop in ArkLib and transfer, is not yet proved; ArkLib components are ported to the local
definitions instead.

\subsection{Perfect completeness}\label{subsec:sec-completeness}

A \emph{strategy} is a deterministic honest prover (no loss of generality against unbounded
adversaries; no randomness is needed without zero knowledge). A \emph{pure verifier} maps statement
and full transcript to an optional output statement, \lean{none} meaning reject. Let
$\mathrm{tr}_\rho$ be the honest transcript on challenge vector $\rho$.

\begin{definition}[Perfect completeness]\label{def:sec-complete}
For every $(s,w)\in\R_{\mathrm{in}}$ and \emph{every} challenge vector $\rho$,
$\Vrf(s,\mathrm{tr}_\rho)=\mathsf{some}\ s'$ where $(s',w')=\Prv.\mathsf{output}(s,w,\mathrm{tr}_\rho)$,
and $(s',w')\in\R_{\mathrm{out}}$.
\end{definition}

\noindent For a deterministic prover, \enquote{for all $\rho$} equals probability one under
uniform challenges, hence ArkLib's \lean{perfectCompleteness}. Bounded-error completeness is not
offered: no IOP-model design needs it (grinding appears only after Fiat--Shamir), and it would bring
probability into the trusted text.

\begin{leancode}[fontsize=\footnotesize,label={\footnotesize Contract/Security.lean (abridged)}]
/-- **Perfect completeness**: for every `(s, w) ∈ relIn` and every challenge vector, the verifier
accepts the honest transcript with the prover's output statement, and the output pair is in
`relOut`. -/
def PerfectComplete (relIn : Set (Stmt × Wit)) (relOut : Set (StmtOut × WitOut))
    (P : Strategy sh Stmt Wit StmtOut WitOut) (V : PureVerifier sh Stmt StmtOut) : Prop :=
  ∀ s w, (s, w) ∈ relIn → ∀ ρ : Challenges sh,
    V s (P.transcript s w ρ) = some (P.output s w (P.transcript s w ρ)).1 ∧
      P.output s w (P.transcript s w ρ) ∈ relOut
\end{leancode}

\subsection{Round-by-round knowledge soundness: the one variant}\label{subsec:sec-rbrks}

\begin{decision}{S-1}{Worst-case RBR-KS with a per-round backward extractor, in counting form}
The contract has exactly one knowledge-soundness notion. It is WARP's definition
\cite[Defs.~4.1--4.2, informally 2.6--2.7]{BCFW25-WARP} (\cref{def:bg-rbrks-warp}) in ArkLib's generalization: a
knowledge state function indexed by an intermediate witness, a per-move backward extractor that is
\emph{named} data of the submission, worst case over statements and prefixes, stated as \enquote{at
most an $\eps_i$-fraction of challenges is bad}.
\end{decision}

\paragraph{Mathematics.} Index the $n$ moves of the protocol by $m\in[n]$ and the stages (prefix
lengths) by $\{0,\dots,n\}$. A \emph{knowledge state function} is a family of predicates
$S_m(s,\mathrm{tr},w)$, for statements $s$, prefixes $\mathrm{tr}$ of length $m$ and intermediate
witnesses $w\in W_m$, with $W_0$ the witness type of $\R_{\mathrm{in}}$. An \emph{extractor} is a
family of maps $E_m: W_{m+1}\to W_m$ (given $s$ and the prefix of length $m+1$) and
$E_{\mathrm{out}}:\ \mathit{WitOut}\to W_n$ (given $s$ and the full transcript). Three laws hold:
\begin{align*}
\textsf{(empty)}\quad & (s,w)\in\R_{\mathrm{in}} \iff S_0(s,\emptyset,w);\\
\textsf{(next)}\quad & \text{if move $m$ is a prover message $a$:}\quad
  S_{m+1}(s,\mathrm{tr}\|a,w)\ \Rightarrow\ S_m\bigl(s,\mathrm{tr},E_m(s,\mathrm{tr}\|a,w)\bigr);\\
\textsf{(full)}\quad & \Vrf(s,\mathrm{tr})=\mathsf{some}\ s' \ \wedge\ (s',w')\in\R_{\mathrm{out}}
  \ \Rightarrow\ S_n\bigl(s,\mathrm{tr},E_{\mathrm{out}}(s,\mathrm{tr},w')\bigr).
\end{align*}
For a challenge move $m$ with prefix $\mathrm{tr}$, a challenge $c\in C_m$ is \emph{bad} if
\[ \exists w\in W_{m+1}:\quad \neg S_m\bigl(s,\mathrm{tr},E_m(s,\mathrm{tr}\|c,w)\bigr)\ \wedge\
   S_{m+1}(s,\mathrm{tr}\|c,w). \]
The verifier is RBR-KS with errors $(\eps_m)$ if
\[ \forall s,\ \forall\ \text{challenge moves } m,\ \forall\ \mathrm{tr}:\qquad
   \#\{c\in C_m : c \text{ bad for } (s,\mathrm{tr})\}\ \le\ \eps_m\cdot|C_m| . \]

\paragraph{The laws, read by a cryptographer.} $S_m(s,\mathrm{tr},w)$ reads \enquote{at this
point of the transcript, $w$ is a valid partial witness}. \textsf{(empty)} says that before any
interaction a valid partial witness is exactly a witness of the relation. \textsf{(next)} says that
a prover message cannot \emph{create} knowledge: if some $w$ is valid after the message, the
extractor pulls it back to a valid witness before the message. \textsf{(full)} says that an
accepting transcript is a state of knowledge, witnessed by the extracted output witness. A
challenge is bad if it creates knowledge out of nothing, that is, if it makes some $w$ valid
although the pulled-back witness is invalid before the challenge. Only challenges can do this, and
RBR-KS says that, from every prefix, few of them do. The consequence is straight-line extraction
along the transcript, mechanized as \lean{chain\_extracts}. Compose $E_{\mathrm{out}}$ and every
$E_m$ backwards from the end of an accepting transcript. If no challenge in the transcript is bad
for its own prefix, \textsf{(full)}, then \textsf{(next)} at prover moves and the negation of
badness at challenge moves, keep the state true down to stage~0, and \textsf{(empty)} yields a
witness. This is the deterministic core of WARP's state-restoration theorem
\cite[Thm.~B.4, Claim~B.6]{BCFW25-WARP}; the rest is a union bound over the moves of the state-restoration
tree.

\begin{leancode}[fontsize=\footnotesize,label={\footnotesize Contract/Security.lean (abridged)}]
/-- A round-by-round extractor with intermediate witness types `WitMid` (ArkLib's
`Extractor.RoundByRound`): `extractMid m` maps a witness after round `m` to one before it. -/
structure Extractor (WitMid : Fin (n + 1) → Type) where
  eqIn : WitMid 0 = Wit
  extractMid : (m : Fin n) → Stmt → Transcript sh m.succ → WitMid m.succ → WitMid m.castSucc
  extractOut : Stmt → FullTranscript sh → WitOut → WitMid (Fin.last n)

/-- A **knowledge state function** for `V` and the named extractor `E`. Laws: `empty` (before
round 0 the state is the input relation), `next` (a prover message cannot make the state true
unless it was true before, with the extracted witness), `full` (an accepted full transcript with
an output pair in `relOut` makes the state true with the extracted final witness). -/
structure KSF (relIn : Set (Stmt × Wit)) (relOut : Set (StmtOut × WitOut))
    (V : PureVerifier sh Stmt StmtOut) {WitMid : Fin (n + 1) → Type}
    (E : Extractor sh Stmt Wit WitOut WitMid) where
  toFun : (m : Fin (n + 1)) → Stmt → Transcript sh m → WitMid m → Prop
  empty : ∀ s w, (s, cast E.eqIn w) ∈ relIn ↔ toFun 0 s (Transcript.empty sh) w
  next : ∀ (m : MessageIdx sh) s (tr : Transcript sh m.1.castSucc) (msg : (sh m.1).Ty) w,
    toFun m.1.succ s (tr.concat msg) w → toFun m.1.castSucc s tr (E.extractMid m.1 s (tr.concat
        msg) w)
  full : ∀ s tr sOut wOut, V s tr = some sOut → (sOut, wOut) ∈ relOut →
    toFun (Fin.last n) s tr (E.extractOut s tr wOut)

/-- A fresh challenge `ch` in round `c` is **bad** for prefix `tr` if some witness makes the state
true after `ch` while the extracted witness does not make it true before `ch`. -/
def KSF.Bad (kSF : KSF relIn relOut V E) (c : ChallengeIdx sh) (s : Stmt)
    (tr : Transcript sh c.1.castSucc) (ch : (sh c.1).Ty) : Prop :=
  ∃ w, ¬ kSF.toFun c.1.castSucc s tr (E.extractMid c.1 s (tr.concat ch) w) ∧
    kSF.toFun c.1.succ s (tr.concat ch) w

/-- **Worst-case RBR knowledge soundness with errors `ε`**, counting form: for every statement,
challenge round and prefix, at most an `ε c`-fraction of the round's challenges is bad. (All
challenge types are finite and nonempty, `Library/Field.lean`, so `Nat.card` is the size.) -/
def RbrKS (kSF : KSF relIn relOut V E) (ε : ChallengeIdx sh → ℝ≥0) : Prop :=
  ∀ s (c : ChallengeIdx sh) (tr : Transcript sh c.1.castSucc),
    (Nat.card {ch // kSF.Bad c s tr ch} : ℝ≥0) ≤ ε c * Nat.card (sh c.1).Ty

/-- The security clause shared by protocols and reductions: perfect completeness, an RBR-KS
state function for the named extractor at errors `ε`, and the error cap `ε c ≤ 2^-128`. -/
structure SecureWith (relIn : Set (Stmt × Wit)) (relOut : Set (StmtOut × WitOut))
    (P : Strategy sh Stmt Wit StmtOut WitOut) (V : PureVerifier sh Stmt StmtOut)
    (E : Extractor sh Stmt Wit WitOut WitMid) (ε : ChallengeIdx sh → ℝ≥0) : Prop where
  complete : PerfectComplete relIn relOut P V
  sound : ∃ kSF : KSF relIn relOut V E, RbrKS kSF ε
  cap : ∀ c, ε c ≤ errorCap
\end{leancode}

\noindent A protocol's security clause instantiates it (\code{Protocol.lean}):
\lean{Protocol.Secure} is \lean{SecureWith (rel sz) Set.univ P.prover P.specVerifier P.extractor P.ε}.
Reading notes: a \lean{Strategy} is a deterministic prover (each message a function of the input
and the prefix; the output a function of the full transcript) and a \lean{PureVerifier} a total
function of statement and full transcript; \lean{Strategy.transcript} is the honest run, built round by round by
\lean{Fin.induction} (\lean{Strategy.prefixes}); \lean{Transcript sh m} is the prefix of \emph{length} $m$, so
\lean{m.castSucc} and \lean{m.succ} are \enquote{before} and \enquote{after} move $m$; \lean{cast
E.eqIn} is an identity coercion; the $\exists w$ in \lean{Bad} ranges over next-stage witnesses;
the state function is existential (proof, not algorithm), the extractor a submission field;
\lean{Nat.card (sh c.1).Ty} is the size of the round's challenge space, finite and non-empty for
every contract round (\code{Library/Challenges.lean}). The same \lean{SecureWith} is used by the
three accumulation reductions (\cref{sec:recursion}).

\paragraph{Why this variant.} The alternatives are those of \cref{rem:bg-variants}.
\begin{itemize}[nosep]
\item \emph{Versus the book's global extractor} \cite[Def.~31.1.6]{CY24-Book} (witness-free state,
  one polynomial-time extractor seeing all prover strings and no challenges): WARP's notion is
  implied by the prefix-extractor notion of CMS19 \cite[App.~C]{BCFW25-WARP}; its relation to the
  book's global extractor is not stated in either source, but WARP's erasure-correction extractor
  needs the next-round witness and so fits only the backward form. It admits codes not known to
  have an efficient error-tolerant decoder, so it excludes no design we want; it is also what
  leanerVM's composition library targets.
\item \emph{Versus ArkLib's averaged and one-shot forms}: the averaged form lets the prover sample
  the prefix, is weaker than the literature's worst case, and puts a prover quantifier (oracle
  simulation, logging) into the audited text. The worst-case form has none, and worst case
  $\Rightarrow$ averaged is proved in ArkLib.
\item \emph{Counting, not probability}: no \lean{Pr}, \lean{ENNReal} or sampling class, and none of
  the \lean{toNNReal}/\lean{toNat} truncation traps. The \lean{Nat.card}${}=0$ trap on infinite
  types is closed because challenge types are the contract's finite, non-empty menu.
\item \emph{No \enquote{iff} in \textsf{(full)}}: WARP states an equivalence, but its proof, checked
  and mechanized by the definitions review (\lean{chain\_extracts}), uses only the directions above;
  the iff would make honest state functions harder and buys nothing.
\end{itemize}

\subsection{The error cap and its derivation}\label{subsec:sec-cap}

\begin{decision}{S-2}{Cap the maximum per-round error at $2^{-128}$; cap rounds at $2^{14}$}
\lean{cap : ∀ c, ε c ≤ errorCap} with \lean{errorCap = (2⁻¹)\^{}128}, a numeral fixed by the
contract, together with \lean{n\_le : n ≤ roundCap} ($2^{14}$) in every protocol and reduction.
\end{decision}

\noindent Without a contract cap, $\eps$ is chosen by whoever proves the theorem, and a declared
error of $1$ makes RBR-KS hold for a verifier that checks nothing (control C7 below). The cap is on
the \emph{maximum} rather than the sum for two reasons. The compiled bound consumes the maximum,
through the factor $(t+k)\cdot\max_i\eps_i$ (\cref{thm:bg-bcs}). And RBR errors concatenate under
sequential composition, so a cap on the sum would be charged $N$ times by an $N$-node
accumulation tree whose compiled bound is unchanged. The interactive knowledge error is then at most
$\sum_i\eps_i\le 2^{14}\cdot2^{-128}=2^{-114}$ \cite[Claim~31.1.3]{CY24-Book}; this is the
$\log k$-bit loss that per-round accounting always incurs, and it is reported.

Over $E=\mathrm{GF}(2^{192})$ algebraic rounds meet the cap with room (degree-3 sumcheck
$2^{-190.4}$; batching $2^9$ claims over a list of $2^6$, $2^{-177}$; unique-decoding proximity at
block length $2^{25}$, $2^{-169}$); only Johnson-regime proximity is tight ($2^{-131}$ to $2^{-134}$
after folding at rate $1/2$, for $\eta\in[0.011,0.018]$). Grinding is unavailable in the IOP model,
so security comes from field size and query count.

\paragraph{The compiled headline.} Let $t$ be the adversary's number of random-oracle queries,
$\lambda=256$ the hash output length, $n\le2^{14}$ the number of moves (the \lean{n\_le} cap of
S-2 counts moves) and $k\le n/2\le2^{13}$ the number of (message, challenge) pairs.
\begin{itemize}[nosep]
\item \emph{Soundness}: worst-case RBR-KS implies witness-free RBR soundness with the same errors
  (take the state $\exists w,\ S_m(s,\mathrm{tr},w)$; every flip event is a bad event). Chiesa--Yogev
  Thm.~31.2.1 then gives state-restoration soundness $(t+k)\cdot\eps$, and the holographic BCS/COS
  theorem, Thm.~32.8.9, gives the argument's soundness error \cite{CY24-Book}.
\item \emph{Knowledge}: WARP Thm.~B.4 with Construction~B.5 (the backward-chain state-restoration
  extractor) \cite{BCFW25-WARP}, and the holographic straight-line knowledge theorem, Chiesa--Yogev
  Thm.~32.8.10, whose hypothesis (Def.~32.6.4) admits an extractor that reads the transcript and
  challenges \cite{CY24-Book}. The extraction time is
  $\sum_i\mathit{et}_i+\mathit{et}_{\mathrm{MT}}+O(r\cdot k\cdot t\log t)$, which is
  \emph{unbounded} at level L1 because $\mathit{et}_i$ is (\cref{subsec:sec-extractor}).
\end{itemize}
The simplified hash term requires $t\ge 2\ell_0$ and $t\ge 2(\log\ell+1)\ell$, where
$\ell=\ell_0+\sum_{j\ge1}\ell_j$ is the \emph{total} length of the index oracle and all prover
oracles \cite[Thm.~32.8.9]{CY24-Book}. Set $t^*:=\max(2\ell_0,\,2(\log\ell+1)\ell)$. The Merkle
and collision terms are non-decreasing in $t$, so for $t<t^*$ they are bounded by their value at
$t^*$; only the quadratic term takes $\max(t,t^*)$, while the first term keeps $t$. Under the
admission caps ($V_k\le2^8$ and the index-oracle cap $\ell_0=\mathit{idxLen}\le2^{k+8}$) this gives
$t^*<2^{43}$ at the top class. The holographic term is included. The headline, stated with these conditions
wherever it appears, is
\[ \boxed{\ \kappa(t)\ \le\ (t+k)\cdot 2^{-128}\;+\;\frac{4.5\,\max(t,t^*)^2}{2^{256}}
   \ \le\ \frac{t}{2^{127}}\qquad\text{for } 2k\le t\le 2^{124}\ } \]
\noindent with $\kappa$ read as a soundness error unconditionally and as a knowledge error with
unbounded extraction time at L1. The earlier draft had \enquote{$k\le t$}, which is false at $k=t$
since the hash term is positive; $2k\le t$ leaves room for it.

\noindent The upper end of the range is where the quadratic term starts to matter: with $k\le t/2$
the bound is at most $1.5\,t\cdot2^{-128}+4.5\,t^2\cdot2^{-256}$, which is at most $t\cdot2^{-127}$
iff $9t\le2^{128}$, i.e.\ $t\le2^{124.8}$; the range is therefore stated as $t\le2^{124}$.

\subsection{The extractor computability problem and honest labelling}\label{subsec:sec-extractor}

\paragraph{The problem.} Lean extractors are arbitrary functions. An extractor that applies
\lean{Classical.choose} to the statement \enquote{a witness exists} collapses knowledge soundness
to soundness. On finite witness spaces a \emph{computable} brute-force extractor does the same:
$K^{n_w}$ is finite and R1CS is decidable, so enumeration is a legitimate Lean function. The
computability gate therefore excludes nothing. The definitions review made this exact.

\begin{theorem}[L1 knowledge equals RBR soundness; A2-M2]\label[theorem]{thm:sec-collapse}
Let $\Vrf$ be any pure verifier with a witness-free RBR state function at errors $\eps$ (counting
form), and let $\mathit{bf}$ be any function that returns a witness whenever one exists. Then
$\mathit{bf}$, used at every stage, is an extractor for which the contract's \lean{RbrKS} holds
with the same $\eps$. Conversely (argued, not mechanised: take $\exists w,\,S_m$), every \lean{KSF}
yields an RBR state function with the same errors.
\end{theorem}

\begin{leancode}[fontsize=\footnotesize,label={\footnotesize Library/RbrSound.lean (untrusted, kernel-checked; abridged)}]
-- `SF`: a witness-free state function (CY Def. 31.1.1 shape); `SF.Rbr ε`: at most ε c · |C|
-- challenges flip the state.
variable (bf : Stmt → Wit) (hbf : ∀ s, (∃ w, (s, w) ∈ relIn) → (s, bf s) ∈ relIn)

/-- **RBR soundness (+ brute force) ⇒ the contract's RBR-KS, same `ε`** (A2-M2). -/
theorem SF.rbrKS_of_rbr (S : SF relIn V) (ε : ChallengeIdx sh → ℝ≥0) (h : S.Rbr ε) :
    RbrKS (S.toKSF bf hbf) ε

/-- **Straight-line extraction** from an accepting transcript with no bad challenge. -/
theorem chain_extracts (kSF : KSF relIn relOut V E) (s : Stmt) (tr : FullTranscript sh)
    (sOut : StmtOut) (wOut : WitOut) (hV : V s tr = some sOut) (hrel : (sOut, wOut) ∈ relOut)
    (hgood : ∀ c : ChallengeIdx sh, ¬ kSF.Bad c s (pre tr c.1.castSucc) (tr c.1)) :
    (s, cast E.eqIn (chain E s tr wOut 0)) ∈ relIn
\end{leancode}

\noindent The collapse is a \emph{weakening, not a vacuity}: unsatisfiable instances exist at
every class, so RBR soundness is a non-trivial property (accept-all has no state function, control
C1), and it is what deployed STARKs are analysed for. Extractor efficiency is load-bearing in two
places: the compiled \emph{knowledge} claim, whose extractor runs the IOP extractor, and recursion,
whose proof-carrying-data bound contains $N\cdot t_E$ for $N$ nodes \cite{CGSY24}. A verifier that
accepts everything on the benchmark index would need a Lean proof that every public input is
satisfiable, for a hash chain a surjectivity proof of the hash; hence the benchmark rule that no
index may be provably always-satisfiable (\cref{sec:benchmark}).

\begin{decision}{S-3}{Label L1 records as RBR-sound; knowledge labels require the L2 track}
Records certified at L1 are labelled \enquote{RBR-sound (knowledge with unbounded extraction)}.
The labels \enquote{knowledge-sound} and \enquote{recursion-ready} require the \emph{L2
efficient-extraction track}: the extractor is also submitted as a RISC-V image, proved to agree
with the Lean extractor and to run within a contract-fixed polynomial cycle bound in the
transcript size. The track is scored on the same leaderboard, with a flag.
\end{decision}

\noindent Rejected, because they shrink the design space: a fixed extractor form (witness in the
clear in a designated oracle range) kills every non-systematic witness encoding, including
FRI-, STIR- and WHIR-style commitments; a fixed decoder for a declared code forces unique decoding.
Rewinding is unnecessary, since extraction that reads the oracles is straight-line.

\subsection{Fiat--Shamir hygiene and sampling}\label{subsec:sec-fs}

The Lean statement quantifies the instance before the prover acts; in the state-restoration game
the adversary chooses $(i,x)$ adaptively \cite[Defs.~32.4.2, 32.6.4]{CY24-Book},
\cite[Def.~B.1]{BCFW25-WARP}, which the worst case over $s$ covers. The compiler must preserve this
binding. Two attack classes sit outside it. Weak Fiat--Shamir and Fenzi's generated-statement
attacks \cite{BPW12-WeakFS,Fen26-GenR1CS} succeed when an accepting transcript can be prepared
before the statement is bound; Fenzi proves that deriving the first challenge from the generated
statement restores the soundness of the underlying protocol. The attack of \cite{KRS25} is
different: it breaks the \emph{strong} transform, with the statement absorbed, whenever the proved
relation can evaluate the Fiat--Shamir hash. No prefix rule prevents it, so it is covered only by
the deployment assumption \lean{RecursionHeuristic} (\cref{subsec:rec-cannot}).

\begin{decision}{S-4}{Fiat--Shamir prefix and bias-free sampling (harness rule)}
Before any challenge of any reduction, the transcript is initialized with
\[ H\bigl(\mathit{domSep}\,\|\,\mathit{contractId}\,\|\,\mathit{shapeHash}\,\|\,\mathit{claimHash}\,\|\,
   k\,\|\,\mathrm{MTroot}(\mathsf{IDX}\,i)\,\|\,\mathsf{VK}\,i\,\|\,x\bigr), \]
and every prover-message root is absorbed before the next challenge, including depth-0 rounds that
are read in the clear. In the accumulation track the full serialized input statement of each
\lean{embed}, \lean{fold} and \lean{decide} run is absorbed likewise. A challenge in $E$ is any
24 hash bytes; an index challenge in $\{0,\dots,2^b-1\}$ is the low $b$ bits of 8 hash bytes.
\end{decision}

\noindent The index root and key bind the circuit, $x$ the instance; shape hash, claim hash and
contract id separate protocols, classes and contract versions; absorbing every root makes each
challenge a function of $(i,x,\Pi_{\le j})$, the hypothesis of the state-restoration theorems. The
Lean theorem assumes exactly uniform challenges, and over binary fields the sampler is exact (every
24-byte string is in $E$; $2^b$ divides $2^{64}$), whereas reducing a 64-bit chunk modulo a 31-bit
prime has a statistical distance of up to $2^{-35}$ per element. Prefix rule and sampler are harness code in the
trusted surface (\cref{tab:arch-trusted}); hash-specific attacks, Merkle and transcript code, and
the hash inside the proved relation are deployment assumptions, never theorem hypotheses.

\subsection{Trusted mathematical hypotheses}\label{subsec:sec-hypotheses}

Lean proofs force proven proximity regimes. Today the unique-decoding regime is proved in ArkLib
without \lean{sorry} (correlated agreement on affine lines \cite{BCIKS20}), as is Johnson-range
mutual correlated agreement for lines. The general Johnson-regime correlated agreement of
\cite{BCIKS20,BCHKS25} and mutual correlated agreement up to Johnson for polynomial generators
\cite{BCGM25-MCA} are published theorems that are not formalized. The line result alone does not
open an unconditional Johnson-regime track: WHIR- and FRI-style folding combine many words with a
polynomial (curve) generator and need mutual correlated agreement for that generator, with explicit
constants, which is the unformalised part. Several \enquote{up to capacity}
conjectures have been refuted in some regimes \cite{DG25-Dist,CS25-PGC,BCHKS25,KKH26,Kam26}. The
proven regime is not free: in the cost model the unique-decoding design pays about $1.7\times$ the
Johnson-regime design at rate $1/2$ (\cref{sec:cost}).

\begin{decision}{S-5}{Unconditional and conditional tracks; capacity conjectures banned}
\begin{itemize}[nosep,leftmargin=1em]
\item \emph{Unconditional track}: kernel-proven facts only, with the standard axioms.
\item \emph{Conditional track}: the contract carries \code{Hypotheses.lean}, a closed, versioned
  list of uninhabited structures. Each is the verbatim transcription of one published theorem
  with a complete proof over the exact code family, with explicit constants and a \emph{regime
  guard} as a field, so that it cannot be instantiated outside its proven range. The first,
  \lean{McaJohnson}, is leanerVM's transcription of \cite[Thm.~4.6]{BCHKS25} (a preprint whose MCA
  proof is given as a sketch for the separable case; the inseparable case, relevant in
  characteristic~2, is argued for the line version by reference to BCIKS20 App.~C), with guard
  $\gamma\le1-\sqrt\rho-\eta$, $\eta>0$, $m=\max(\lceil\sqrt\rho/\eta\rceil,3)$, the conservative
  Thm.~4.6 constant. Its admission is an organiser exception, recorded as such, until a refereed
  version or a complete characteristic-2 proof exists; the refereed alternative is BCGM25's MCA up
  to Johnson \cite{BCGM25-MCA} (CCC 2026, any finite field), whose exceptional count is quadratic
  in $n$ and costs bits (\cref{tab:bg-binary-errors} is not recomputed under it). The conditional stub is
  \lean{theorem certificate (h : Hypotheses) : Certificate …}, and records are flagged
  \enquote{conditional on H-v$n$}.
\item \emph{Banned in both tracks}: any capacity or list-decoding-capacity conjecture.
\end{itemize}
\end{decision}


\noindent \code{Hypotheses.lean} is planned, not part of contract v0.3: \lean{McaJohnson} will be an
uninhabited \lean{Prop} structure bounding \lean{mcaError} of RS codes over $E$ on subspaces of $K$
under the guard, and the conditional track opens only when it is added under the two-reviewer rule. The guard is load-bearing: for every characteristic-2 field there are Reed--Solomon
codes on essentially random $\mathbb{F}_2$-subspaces (at rate $\approx1/16$, Cor.~1.7) with
$n^{2-o(1)}$ exceptional points \emph{at} the Johnson radius \cite[Thm.~1.6, Cor.~1.7]{BCHKS25}, so
no bound linear in $n$ holds at $\eta=0$ for all such codes; with $\eta>0$ the $O(n/\eta^5)$ bound applies. Governance: a new hypothesis needs two
organizer reviewers, a published proof and a check against the known negative results (no tool
detects an inconsistent hypothesis). A hypothesis later formalized is inhabited in the library and
its records become unconditional on re-check; a refuted one is removed and its records revalidated
(\cref{subsec:arch-versioning}). \enquote{Everything proven} was rejected: it penalizes designs
resting on published but unformalized theorems and pushes agents to unique-decoding query counts
($1.2$--$2.2\times$ the Johnson counts, \cref{tab:bg-queries}) or to
enormous kernel-evaluated certificates.

\subsection{Non-vacuity}\label{subsec:sec-nonvacuity}

Kernel checking certifies the statement, not its meaning (\cref{subsec:bg-lean}). Every clause of
the certificate was therefore attacked during design review, and \cref{tab:sec-vacuity} lists how
each could be vacuous or weaker than it reads, and what closes it. Entries marked A2 were Lean
counterexamples against an earlier draft; each is now a control that must fail against the
contract.

\begingroup\small
\renewcommand{\arraystretch}{1.12}
\begin{longtable}{@{}>{\raggedright\arraybackslash}p{0.17\textwidth}>{\raggedright\arraybackslash}p{0.41\textwidth}>{\raggedright\arraybackslash}p{0.36\textwidth}@{}}
\caption{How each statement could be vacuous, and what closes it.}\label{tab:sec-vacuity}\\
\toprule
Statement & How it could be vacuous or weak & Closed by \\
\midrule\endfirsthead
\toprule
Statement & How it could be vacuous or weak & Closed by \\
\midrule\endhead
Completeness & empty relation at some class; empty challenge type (no run exists) & relation non-empty at every class (zero index); challenge types are contract types, non-empty \\
Completeness vs.\ soundness & reject-all is sound; accept-all is complete & both are required; controls C1, C3 \\
RBR-KS & $\eps=1$; \lean{Nat.card}$=0$ on an infinite type; $\eps\cdot0$ on an empty type & numeral cap $2^{-128}$; finite, non-empty contract challenge types; C7 \\
Round count & unbounded rounds, each at the cap, accumulate to soundness error $\ge1/2$ (A2: \lean{hitRed\_secure}) & \lean{n\_le : n ≤ roundCap} ($2^{14}$) in every protocol and reduction; control \lean{hit\_accepts\_lt\_half} \\
Message types & typed messages restrict the adversary; the machine is unconstrained off the image of the serializer (A2: \lean{flagMachine\_accepts\_false}; A1: typed promise) & raw vectors \lean{OVal o}; no serializer (\cref{subsec:proto-rounds}); control \lean{flagged\_attack\_fails} \\
Challenge types & submitter types; out-of-image tapes; free harness work (A2-C2, A1-M1) & contract menu $E^{n_E}\times\prod_b\{0,\dots,2^b-1\}$ with contract encodings \\
Relation & leanerVM: theorems hold of instances whose relation is empty & the relation is the contract's, over every index; unsatisfiable instances at every class \\
Indexer & infeasible key or index oracle (e.g.\ a satisfiability table, which is computable in Lean's sense) & badge-path evaluation under budget (records only); the gate ensures compiled code exists \\
Key size & huge \lean{VK i} as lookup table & \lean{vkLen ≤ 2\^{}10} in \lean{Protocol} \\
Extractor & \lean{Classical.choose} or brute force & not closable at L1: honest labelling (\cref{thm:sec-collapse}); L2 cycle bound \\
\lean{Agrees} & only typed transcripts (see above); spec ignores the transcript; machine never halts & raw quantifier; a spec ignoring the transcript is sound only if it rejects all, which completeness forbids; fuel exhaustion rejects, so an image that loops rejects honest proofs and fails \lean{Agrees} \\
Accepting-run cost & bound vacuous where no run accepts & completeness + \lean{Agrees}: every honest run accepts, so the bound covers it; unsatisfiable indices are irrelevant to cost \\
Accumulator size & cap on a serializer the definitions never read (A2: \lean{trivialAcc0\_secure}) & fixed accumulator shape: plain part $\le 2^{12}$ elements, declared oracles; output agreement (\cref{sec:recursion}) \\
Accumulator oracles & untyped oracle lists: the machine charges the length it sees, so no \lean{Certificate2} exists for a decider that queries an input oracle (A4-F1: \lean{no\_certificate2}) & oracles typed by the declared shape (\lean{AccStmt}); positive control \lean{qimg\_decide\_cost\_clause} \\
Security vs.\ cost & cost clause inside the security certificate & security clauses never mention cost, so security cannot be traded for cost \\
Hypotheses & inconsistent or out-of-regime hypothesis & regime guards; two-reviewer rule; separate, flagged track \\
Scalars and image & score inputs or image bytes read from the solution by compiled evaluation, which \lean{implemented\_by} can make differ from the kernel's value (A4-F2: \lean{UnboundImage}) & image, shape summaries, counters and scores are rendered organiser data; \lean{summary} and \lean{score} clauses prove them; the scorer only takes roots \\
Hygiene & \lean{sorry}, \lean{native\_decide}, \lean{implemented\_by} hidden from a source-text scan; a solution module named like a contract module (A4-F3; \lean{SpoofCert}) & environment gate trusting organiser modules only if their files match the organiser hash manifest, rejecting \code{init}/\code{export} entries; comparator first in a fresh tree; kernel replay (\cref{subsec:arch-ci}; A6-G0--G2) \\
\bottomrule
\end{longtable}
\endgroup

\paragraph{Controls.} The contract repository carries kernel-checked, untrusted controls, run in
continuous integration (\cref{tab:sec-controls}). A \emph{negative} control must \emph{fail} the
certificate (refuted by a proved theorem, or ill-typed against the contract); a \emph{positive}
control must satisfy it, showing that the certificate is inhabited and means what it says.

\begin{table}[h]
\centering\small
\begin{tabularx}{\linewidth}{@{}l>{\raggedright\arraybackslash}X@{}}
\toprule
& Negative controls (must fail) \\
\midrule
C1 & accept-all verifier: no state function on one-message schedules (\lean{acceptAll\_no\_KSF}); error $\ge1$ at the first challenge otherwise (\lean{acceptAllChal\_err\_ge\_one}) \\
C2 & an unsatisfiable instance at every class (\lean{badIndex\_unsat}: $1\cdot1=0$ on the constant wire) \\
C3 & reject-all verifier is not complete (\lean{rejectAll\_not\_complete}) \\
C6 & a verifier skipping the last constraint row fails RBR-KS (planned; not yet in v0.3) \\
C7 & cap is load-bearing: with cap 1, accept-all on a one-challenge schedule \emph{is} RBR-KS (\lean{acceptAllChal\_rbr\_one}, \lean{acceptAllChal\_not\_secure}, \lean{errorCap\_lt\_one}) \\
C9 & machine: a looping image, or one trapping on an honest transcript, fails \lean{Agrees} (planned; the machine rejects of \lean{MachineSmoke} cover HALT(0), HASH, a missing oracle and a store into the tape) \\
A2-A/B/C & flag-word attack fails against raw messages (\lean{flagged\_attack\_fails}); hitting protocol: secure at error $2^{-192}$ (\lean{hit\_secure}), yet with $n\le2^{14}$ rounds it accepts the false statement on fewer than half of all transcripts (\lean{hit\_accepts\_lt\_half}); accumulator outputs only selected oracles (\lean{verifier\_oracles\_selected}, \lean{empty\_acc\_output}) \\
H & harts: each race or discipline violation rejects (\lean{child\_escapes}, \lean{overlap\_rejected}, \lean{early\_read}, \lean{unjoined}, \lean{no\_flag}) \\
G & A4 probes: \lean{GateBypass}, \lean{UnboundImage}, \lean{HygieneBypass}, \lean{SpoofCert} rejected by the gate; \lean{SkipKernel} rejected by \code{leanchecker} in CI. A6 regressions rejected: \code{cmpproj} (comparator order) and \code{gateproj} (\lean{InitSol}, \lean{ExportSol}, \lean{TamperSol}/\lean{TamperCert}); 22 malformed claims (13 schema, 7 path, 2 A6) rejected by the renderer \\
\midrule
& Positive controls (must hold) \\
\midrule
C4 & toy \enquote{send the witness} protocol is secure over the sparse index (\lean{toy\_secure}, \lean{toyTest\_secure}); the certified baseline (security level L0, \cref{sec:baseline}) is the machine part \\
C5 & \lean{fullRecAcc} (embed = whole verifier, fold = decide = accept-all) satisfies \lean{Accumulation.Secure} (\lean{fullRecAcc\_secure}): the T2 baseline \\
C8 & bridge: local $\Rightarrow$ ArkLib worst-case RBR-KS (\lean{RbrKS.toArk}); the converse and a completeness bridge are open \\
E, F & \lean{chain\_extracts} and its corollary (knowledge error $\sum_i\eps_i$); \lean{K\_spec}, field cardinalities; \lean{pad\_sat\_iff} \\
M & machine, by kernel evaluation at full fuel: \lean{smoke\_accepts} (18 counters), \lean{smoke\_timing}, \lean{firsttouch\_counts}, \lean{firsttouch\_fewer}; harts \lean{disjoint\_accepts}, \lean{span\_lt\_work} \\
F1 & an oracle-reading decider over a typed accumulator meets an exact cost bound on every input (\lean{qimg\_decide\_cost\_clause}) \\
\bottomrule
\end{tabularx}
\caption{Negative and positive controls.}
\label{tab:sec-controls}
\end{table}
